\begin{abstract}
High-capacity color QR codes stack three
binary layers in the RGB channels, and decoding them depends on recovering each module's
color from a phone photograph. We present a color recovery method that needs no offline
training. It inverts a one-parameter linear leak between neighboring modules and then
classifies colors with regularized discriminant analysis, fitted per photograph on the
function patterns whose colors the symbol fixes. Both parameters were chosen on each of two
datasets separately, and the chosen values are best or statistically tied with the best on
both: the public CUHK-CQRC dataset of HiQ codes and a new dataset of \QrAttempted{}
photographs, which we release. On CUHK-CQRC the method decodes
\HiqCodeChosen\% of localized symbols. On HiQ's own test split, with HiQ's
shipped classifiers run on the same module samples, ours decodes \PortOursCode\% of symbols
against \PortCmiCode\% by HiQ's best, at \TimeChosenNs{} against \PortCmiNs\,ns per module on
one Apple M4 core, our per-photograph fit included (\TimeChosenVforty\,ms for a version-40
symbol). End to end, our open-source decoder reads
\PairCommonOurs\% of the symbols that both it and the released HiQ decoder localize, against
\PairCommonHiq\% for HiQ.
\end{abstract}

\begin{IEEEkeywords}
color QR code, high-capacity barcode, color recovery, discriminant analysis, mobile imaging
\end{IEEEkeywords}
